\documentclass[aspectratio=169, 10pt] {beamer}
\usetheme[faculty=phil]{fibeamer}
\usepackage[utf8]{inputenc}
\usepackage[main=english,]{babel}
\usepackage{subcaption}

\makeatletter
\renewcommand\fibeamer@includeLogo[2][]
\makeatother

\title{A Tool for Signal Chain Simulation in Particle Physics}
\subtitle{Octavio Gamez and Desiree Martinez} 
\author{July 31st, 2026}

\usepackage{ragged2e}  % `\justifying` text
\usepackage{booktabs}  % Tables
\usepackage{tabularx}
\usepackage{tikz}      % Diagrams
\usetikzlibrary{calc, shapes, backgrounds}
\usepackage{amsmath, amssymb}
\usepackage{url}       % `\url`s
\usepackage{listings}  % Code listings
\frenchspacing
\includegraphics[trim=left bottom right top, clip]{filename}

\begin{document}
\frame[c]{\maketitle}

%Slide 2
    \begin{frame}{Motivation}
    \large
    Experiments focus on the signal to the digitizer, but hardly consider how the signal travels and what happens to the signal\\

    %ADD PICTURE OF LIKE A PMT TO DIGITIZER TYPA TOOL IDK
    \end{frame}


%Slide 3
    \begin{frame}{Gaps in Pre-Existing Tools}
    SPICE simulators are accurate but not ideal or realistic during the early planning stages of an experiment\\

    %ADD PICTURE OF A SPICE SIMULATOR 

    \begin{figure}
        \hfill
        \includegraphics[width=0.7\linewidth]{SPICE.jpg}
        \caption{Caption}
        \label{fig  1:placeholder}
    \end{figure}


    \small
    Figure 1: Electrical Engineering Oli Glaser
%I cant figure out how to show this, i need the pic to be big but my text wont freaking go to the left side of the screen
    
    %DURING PRESENTING WE CAN TALK ABOUT HOW IT IS ACCURACT BUT REQUIRES A FULLY DETAILED SCHEMATIC AND YEAH YA FEEL?
    %MAYBE GO ONTO HOW THIS IS ONE OF THE OPTIONS BUT THE OTHER ONE IS TO JUST KEEP GOING WITH YOUR GUT AND  PAST EXPERIENCES WHICH ISNT IDEAL
    \end{frame}


%Slide 4
    \begin{frame}{Project Goal}
    \large
    Design an open source tool for the pre-planning stage of an experiment that will report reflections, distortions, and signal loss

    %IDK WHAT PICTURE YOUD WANT TO ADD
    \end{frame}


%Slide 5 (MAYBE COULD ADD IDK IF IT WOULD BE RELEVENT BUT IT DONT MATTER
%    \begin{frame}{Understanding a signal}
%    Experimenters try to understand things they see in signals; 
%    \begin{itemize}
%        \item Source
%        \item Timing
%        \item Reflections
%        \item Distortions
%        \item Signal Loss
%    \end{itemize}

    %ADD PIC OF SIGNAL AND POINT OUT THE PARTS OF THE SIGNAL IDK
%    \end{frame}
%GET RID OF THE PERCENTAGE SIGNS IF YOU WANT TO ADD IT



\begin{frame}{PMT Signal}
$$
f(t) = A \cdot \left[ e^{-\frac{t}{\tau_{\text{fall}}}} - e^{-\frac{t}{\tau_{\text{rise}}}} \right]
$$
Where:
\begin{align*}
f(t) &\quad\to \text{output values at a specific time } (t) \\
A    &\quad\to \text{an arbitrary scaling factor} \\
t    &\quad\to \text{time} \\
\tau_{\text{rise}} &\quad\to \text{Time constant controlling the signal rise time}\\
\tau_{\text{fall}} &\quad\to \text{Time constant controlling the signal fall time}
\end{align*}

% add a picture of what a PMT signal looks like
% maybe can talk about the fast rise with a slow fall time along with a negative going pulse
        %like we can actually talk about it or we can put in on the slide and still talk about it
\end{frame}



%Slide 6
\begin{frame}{Impedance}
Impedance Equation:
$$Z_0 = \sqrt{\frac{R + j\omega L}{G + j\omega C}}$$
\end{frame}



%Slide 7
\begin{frame}{Reflections}
    Reflection Coefficient:
$$\Gamma = \frac{(Z_{\text{load}} - Z_0)}{(Z_{\text{load}} + Z_0)}$$

 A matched connection ($Z_0$ = $Z_{\text{load}}$) gives no reflection ($\Gamma$ = 0). An open or short end cable yields:

$$\Gamma_{\text{open}} = \lim_{Z_{\text{load}} \to \infty} \frac{Z_{\text{load}} - Z_0}{Z_{\text{load}} + Z_0} = +1$$

for an open cable, and: 

$$\Gamma_{\text{short}} =  \frac{0- Z_0}{0 + Z_0} = -1$$
\end{frame}




%Slide 8
\begin{frame}{Transmission Line}
Propagation delay is given by:

$$t_{delay} = \frac{L}{v_f \cdot c}$$

Where:
\begin{align*}
t_{delay} &\quad\to \text{the signal's shift in seconds} \\
L   &\quad\to \text{length of the cable in meters} \\
v_f    &\quad\to \text{velocity factor (ratio of transmit speed to the speed of light)} \\
c    &\quad\to \text{the speed of light)} \\
\end{align*}
\end{frame}



\begin{frame}{PMT}
    PMT &\quad\\  

    \begin{figure}
        \centering
        \includegraphics[trim=3cm 2cm 1.9cm 3cm, clip, width=1.0\linewidth]{PMTSignal.png} 
        \caption{Synthetic PMT Pulse}
        \label{fig:placeholder}
    \end{figure}
    
\end{frame}




\begin{frame}{Cable}
    PMT \to \text{1M 50 $\Omega$ cable} \\
     \begin{figure}
        \centering
        \includegraphics[trim=3cm 2cm 1.9cm 3cm, clip, width=1.0\linewidth]{AMP_IN.png} 
        \caption{Synthetic PMT Pulse}
        \label{fig:placeholder}
    \end{figure}
\end{frame}


\begin{frame}{Amplifier}
    PMT \to \text{1M 50 $\Omega$ cable} \to Amplifier \\
     \begin{figure}
        \centering
        \includegraphics[trim=3cm 2cm 1.9cm 3cm, clip, width=1.0\linewidth]{AMP_OPEN_CIRCUIT.png} 
        \caption{Synthetic PMT Pulse}
        \label{fig:placeholder}
    \end{figure}
\end{frame}


\begin{frame}{Cable}
    PMT \to \text{50 $\Omega$ cable} \to Amplifier \to \text{5M 50 $\Omega$ cable} \\
     \begin{figure}
        \centering
        \includegraphics[trim=3cm 2cm 1.9cm 3cm, clip, width=1.0\linewidth]{DIG_ANALOG_INPUT.png} 
        \caption{Synthetic PMT Pulse}
        \label{fig:placeholder}
    \end{figure}
\end{frame}



\begin{frame}{Next Steps}
\begin{itemize}
    \item Test the tool against real measurements to check accuracy and compare
    \item Get it to support multiple reflections bouncing back and fourth
    \item Add more components and source types
    \item enhancing different features that are already implemented in our design
\end{itemize}
    
\end{frame}


\begin{frame}{Conclusion}
    \begin{itemize}
        \item 2 paths when planning an analog signal:
        \begin{itemize}
            \item experience-based estimations
            \item building a fully designed simulation
        \end{itemize}
        \item This project was done to add another path
        \begin{itemize}
            \item open-source tool
            \item has a source and components that could distort a signal
            \item giving back results before hardware is built
            \begin{itemize}
                \item reflections, distortion, signal loss
            \end{itemize}
        \end{itemize}
    \end{itemize}
    
\end{frame}




\end{document}
